% ECONOMICS_PROBLEM: {"id": "D13-434", "title": "Measuring inequality within households", "jel_code": "D13", "source_id": "OP-0458"}
% !TeX program = xelatex
\documentclass[11pt,a4paper]{article}
\usepackage[margin=23mm]{geometry}
\usepackage{fontspec}
\setmainfont{Latin Modern Roman}
\usepackage{amsmath,amssymb,mathtools}
\usepackage{xcolor,enumitem,booktabs,longtable,array,xurl,hyperref}
\definecolor{muted}{HTML}{53636C}
\hypersetup{colorlinks=true,urlcolor=blue,linkcolor=blue}
\setlength{\parindent}{0pt}
\setlength{\parskip}{5pt}
\newcommand{\R}{\mathbb R}
\newcommand{\E}{\mathbb E}
\newcommand{\N}{\mathbb N}
\newcommand{\ind}{\mathbf 1}
\newcommand{\Var}{\operatorname{Var}}
\newcommand{\Cov}{\operatorname{Cov}}
\newcommand{\supp}{\operatorname{supp}}
\DeclareMathOperator*{\argmin}{arg\,min}
\DeclareMathOperator*{\argmax}{arg\,max}
\newcommand{\entrymeta}[3]{{\sffamily\small\color{muted}\textbf{\texttt{#1}}\quad|\quad #2\par #3\par}}
\newcommand{\Problemref}[1]{\texttt{#1}}
\newcommand{\JELref}[2]{\texttt{#2} (JEL \texttt{#1})}
\begin{document}
\section*{Measuring inequality within households}
\textbf{Problem ID:} \texttt{D13-434}\quad \textbf{JEL:} \texttt{D13}\par
% BEGIN ECONOMICS STATEMENT
\label{problem:OP-0458}
{\sffamily\small\color{muted}Original subject group: Education, families and demography\par}
\label{definitions:problem:30007740}
\textbf{Audit classification:} Background; proposed extension. Chapter identity/authors verified through publisher contents; exact individual-poverty passage and DOI not recovered. Defined reference expenditure functions.
\subsection*{Definitions and precise research target}
\textit{Editorial formalization.} Consider two-adult households with children in one nationally representative expenditure survey. Let household \(i\) have disposable resources \(R_i\), observed assignable adult expenditures \(x_{i1},x_{i2}\), and joint expenditures \(g_i\). Let latent individual welfare be \(u_{ij}(c_{ij},g_i)\), where \(c_{ij}\) is adult \(j\)'s private consumption and \(j\in\{1,2\}\). For an individually defined poverty threshold \(z>0\), let \(m_{ij}\) be the monetary equivalent of this welfare under a specified common reference price vector. The target is
\[P=E\!\left[\frac{\mathbf1\{m_{i1}<z\}+\mathbf1\{m_{i2}<z\}}{2}\right].\]
Identify \(P\), or sharp bounds on it, from expenditures, time use, and any available individual reports. Explicitly state restrictions on sharing rules, public-good consumption, preferences, and measurement error; household income alone generally does not determine individual monetary equivalents. Validate inferred allocation against independent individual outcomes, and quantify sensitivity to alternative economies of scale.

The target covers the two adults, not the children's poverty. Let the private consumption bundle \(c_{ij}\in\mathbb R_+^K\), jointly enjoyed service vector \(g_i\in\mathbb R_+^G\), and adult preferences \(u_{ij}\) be declared. A money metric needs an expenditure function \(e_{ij}(p^*,u)=\inf\{p^*\cdot q:u_{ij}^{\rm ref}(q)\geq u\}\), with a specified reference technology translating joint goods to individual services, price vector \(p^*>0\), and reference preferences. Then \(m_{ij}=e_{ij}(p^*,u_{ij}(c_{ij},g_i))\). Assume finite feasible expenditures; ordinal utility alone does not choose the common reference environment. Adult welfare may differ despite identical household resources. Sharp bounds mean the infimum and supremum of \(P\) across all preference/allocation/measurement models compatible with the observed law and stated restrictions; no unidentified sharing rule is silently fixed. All expectations use a stated baseline population and probability law. Identification means every admissible data-generating model with the same observed-data law yields the same displayed target; otherwise report the identified set. Exact equations, horizons and assumptions are editorial, rather than source-given canonical conjectures.
\subsection*{Variants and assumption changes}
\begin{enumerate}
\item \textbf{Editorial efficient-allocation variant.} Impose maximization of \(\mu_i u_{i1}+(1-\mu_i)u_{i2}\), \(\mu_i\in(0,1)\), with identified distribution-factor restrictions. Compare poverty bounds with noncooperative allocation.
\item \textbf{Editorial public-good variant.} Compare jointly enjoyed housing/childcare with excludable private goods at common prices. A household equivalence scale is an additional allocation assumption.
\end{enumerate}
\subsection*{References and source audit}
\textbf{1.} Chapter identity and allocation topic verified; exact individual-poverty target is editorial. Locator: Handbook of the Economics of the Family, Volume 1, Chapter 3; publisher table of contents and indexed overview. Chapter full text inaccessible; author list verified on publisher book page. DOI not independently verified and omitted. URL: \url{https://www.sciencedirect.com/science/chapter/handbook/pii/S2949835X23000083}.
\begin{enumerate}\raggedright
\item\hypertarget{definitions:source:30007740:1}{} Ingvild Almås; Orazio Attanasio; Pedro Carneiro. 2023. \emph{Household decisions and intra-household distributions}. Handbook of the Economics of the Family, Volume 1, Chapter 3; publisher table of contents and indexed overview.
\url{https://www.sciencedirect.com/science/chapter/handbook/pii/S2949835X23000083}.
\end{enumerate}

\subsection*{Common conventions: Source mathematical conventions}

These conventions apply to each entry unless explicitly replaced.

\textbf{Models and identification.} A primitive model \(m\) specifies all probability laws, feasible choices, information, equilibrium and measurement rules used by its target. The admissible class \(\mathcal M\) is fixed independently of outcomes. If \(P_O(m)\) is its observable probability law and \(\tau(m)\) the target, its identified set at observable law \(P\) is
\[
 \mathcal I_\tau(P)=\{\tau(m):m\in\mathcal M,\ P_O(m)=P\}.
\]
Point identification means that this set is a singleton. An empty set signals incompatibility of data and assumptions. Coordinate bounds need not describe the joint set. Bounds are sharp only if every retained target is attainable or arbitrarily approachable by compatible models. Finite bounds require explicit restrictions on unobserved quantities; unrestricted confounding, hidden wealth, extrapolation or arbitrary equilibrium selection can yield uninformative sets.

\textbf{Causal interventions.} A potential outcome \(Y_i(p)\) is the outcome under a complete policy regime \(p\), including timing, financing, information and market spillovers. The target population and units are fixed before treatment. With interference, use the complete assignment vector or a justified exposure mapping; individual no-interference notation alone cannot identify an economy-wide intervention. A difference of mean potential outcomes does not require the cross-world joint distribution. Conditioning on post-treatment migration, survival, enrollment or employment changes the estimand and needs separate justification.

\textbf{Statistical statements.} An estimator, confidence set, forecast or test specifies a sampling law, model class and loss/coverage criterion. Uniform \((1-\alpha)\)-coverage means \(\inf_{m\in\mathcal M}\Pr_m\{\tau(m)\in C_n\}\geq1-\alpha\), or an explicitly asymptotic version. The whole parameter space trivially has coverage, so informativeness requires a stated width, risk or power objective. A forecasting improvement is relative to a named benchmark, information set and evaluation population.

\textbf{Equilibrium and optimization.} An equilibrium comprises feasible plans, optimality, beliefs where needed, budget constraints and all market/resource clearing. The primitives or a stated benchmark define each correspondence \(\mathcal E(p;m)\). Nonemptiness is assumed or proved, never inferred just from compact policy sets. A selected-equilibrium target uses a declared selection rule; a robust target quantifies over all permitted equilibria. Use suprema or infima unless attainment is established. An \(\varepsilon\)-optimal policy is feasible and within \(\varepsilon\) of the finite optimum value. Report an empty feasible set separately.

\textbf{Welfare and accounting.} Normative utility, interpersonal normalization, social weights, numeraire, discounting and the use of public receipts are primitives. Behavioral utility need not coincide with the welfare criterion. Household welfare includes ownership income and rents once; adding profits again to compensating variation generally double-counts them. A monetary equivalent is defined by an explicit transfer and price convention, with monotonicity and range conditions. Average effects, mechanism contributions, transfers and welfare losses are different targets. Infinite sums require convergence and dynamic feasible plans require terminal or no-Ponzi conditions.

\textbf{Source versions.} A working-paper date, revision date and journal publication year are distinguished. A DOI for a journal version is not automatically the DOI of an earlier manuscript. Locators refer to the stated version and printed pages where specified. A metadata-only check does not verify a claim about a full-text conclusion. Every entry's source subsection narrows the provenance claim accordingly.
% END ECONOMICS STATEMENT
\end{document}
